\section{Conclusion}
\label{sec:conclusion}

Robot learning is organising itself around a single question: should competence be shipped
as \emph{weights} or as \emph{skills}? This survey took that question as its axis. We mapped
the field into six branches (\S\ref{sec:taxonomy}), positioned end-to-end vision-language-
action models as the weights-based foil (\S\ref{sec:vla}), and gave the code-as-policy camp
its own deep-dive (\S\ref{sec:code}) organised by \emph{degree of self-improvement}: from
one-shot program synthesis, through within-task repair, persistent skill memory, and
population search, to the open-ended loop that combines all three. That top cell, feedback
$+$ memory $+$ search, is the least populated and least surveyed region of the field, and
it is where systems such as ASPIRE, ENPIRE, and RoboClaw now sit. We complemented this with a map of
how skill \emph{repertoires} are built (\S\ref{sec:skill}), showing that skills come in an
RL flavour and a code flavour of which only the latter self-improves without gradients, and
we connected the academic taxonomy to the nascent skill economy (\S\ref{sec:open}), where
commercial marketplaces already distribute skills but ship only static playback. The gap
between that static distribution and genuine on-device adaptation is, we argue, the defining
research opportunity for physical AI, and the self-improvement ladder of \S\ref{sec:code} offers a structured path toward closing it.
